% Part: computability
% Chapter: computability-theory
% Section: computable-sets

\documentclass[../../../include/open-logic-section]{subfiles}

\begin{document}

% SEGMENT OLP-0237-S01

\olfileid{cmp}{thy}{cps}
\olsection{கணிக்கத்தக்க கணங்கள்}

கணிப்புத்தன்மை என்ற கருத்தைக் கணிக்கத்தக்க சார்புகளிலிருந்து
கணிக்கத்தக்க கணங்களுக்கு நீட்டிக்கலாம்:

\begin{defn}
  $S$ இயல் எண்களின் ஒரு கணம் ஆகட்டும். அதன் சிறப்பியல்புச்
  சார்பு~$\Char{S}$ கணிக்கத்தக்கது எனில், எனில் மட்டுமே,
  $S$ \emph{கணிக்கத்தக்கது}. வேறுவிதமாகக் கூறினால், பின்வரும்
  சார்பு கணிக்கத்தக்கது எனில், எனில் மட்டுமே,~$S$ கணிக்கத்தக்கது:
\[
\Char{S}(x) =
\begin{cases}
1 & \text{$x \in S$ எனில்} \\
0 & \text{இல்லையெனில்}
\end{cases}
\]
அதேபோல், உறவு $R(x_0, \dots, x_{k-1})$ அதன் சிறப்பியல்புச்
சார்பு கணிக்கத்தக்கதாக இருந்தால், எனில் மட்டுமே, கணிக்கத்தக்கது.

கணிக்கத்தக்க கணங்களும் உறவுகளும் \emph{தீர்மானிக்கத்தக்கவை}
என்றும் அழைக்கப்படுகின்றன.
\end{defn}

\begin{explain}
இப்போது பகுதிச் சார்புகள், சார்புகள், கணங்கள் ஆகியவற்றுக்கான பல
கணிப்புத்தன்மைக் கருத்துகள் நம்மிடம் உள்ளன என்பதைக் கவனிக்க.
அவற்றைக் குழப்பிக்கொள்ள வேண்டாம்!{} \iftag{TMs}{ஒரு பகுதிச்
சார்பைக் கணிக்கும் டியூரிங் பொறி, சார்பு வரையறுக்கப்பட்டுள்ள
உள்ளீட்டு மதிப்புகளுக்கு அச்சார்பின் வெளியீட்டைத் திருப்பித்தருகிறது;
ஒரு கணத்தைக் கணிக்கும் டியூரிங் பொறி, உள்ளீட்டு மதிப்பு அக்கணத்தில்
உள்ளதா இல்லையா என்று தீர்மானித்த பின் $1$ அல்லது~$0$ ஐத்
திருப்பித்தருகிறது.}{}
\end{explain}

\end{document}
